% GENERATED, edit the French unit or the translation table
\mxunit{seconde-annexes}{Teacher annexes: use and correspondence}
\section*{Using the common source}
\begin{keypoint}[One source, two versions]
The class option \texttt{edition=student} produces the student version (\texttt{eleve-en.tex}); \texttt{edition=teacher} produces the teacher version (\texttt{professeur-en.tex}). The statements and course content are common. Answers, written-out proofs, corrected figures and pedagogical cues appear only in teacher mode.
\end{keypoint}
All graphs and diagrams are written in TikZ/PGFPlots. No external image or page of the original PDF is needed for compilation. Compile twice with PDFLaTeX to update the table of contents and the links. The two PDFs may have different pagination; the stable references are the sequence and example numbers.
\begin{method}[Implementation cues]
\begin{itemize}
\item Before an example: have students restate the data and the question; allow a preparatory representation.
\item During the investigation: provide targeted help (first step, diagram, simple case) and an extension for students who have finished.
\item During the correction: separate result, method, justification and verification. Have a frequent error explained.
\item For fluency work: mix recent and older questions, then redo the fragile ones after a delay.
\item For language: make explicit ``among'', ``at least'', ``if'', ``and'', ``or'' and the meaning of the symbols before calculating.
\end{itemize}
\end{method}
\section*{Common criteria for a piece of work}
\begin{center}\small\mxarraystretch{1.6}\begin{tabularx}{\linewidth}{l X}
Search&An approach is under way; the attempts provide information.\\
Model&The variables, domain and assumptions match the situation.\\
Represent&The table, graph or diagram is readable, labelled and consistent.\\
Reason&Properties are applied under their conditions; the conclusion follows.\\
Calculate&The operations are correct; exact values are kept when necessary.\\
Communicate&The write-up answers the question, with appropriate units and precision.
\end{tabularx}\end{center}
For each criterion, four levels can be used: not started, partial, correct with help, independent. This grid is used to give usable feedback; it does not replace a marking scheme specific to an assessment.

\section*{Correspondence with the original booklet}
\begin{center}\small\mxarraystretch{1.5}\begin{tabularx}{\linewidth}{p{6.3cm} X}
\textbf{Sequences of the supplied booklet}&\textbf{Destination in this adaptation}\\\hline
1. Algebraic calculation&8; reactivation in the starter routines from the start of the year.\\
2. Numbers and calculations&1 for arithmetic; 4 for sets; 8 for powers, fractions and roots.\\
3. Intervals and absolute value&4; reuse of the absolute value function in 6 and 10.\\
4. Equations and inequalities&8; simple equations reactivated in 3 and 6.\\
5. Proportion and percentage change&9.\\
6. Vectors&2.\\
7. Functions and graphical representation&6 for representations and readings; 10 for variations and extrema.\\
8. Notions of probability&14; set vocabulary introduced in 1.\\
9. Notions of collinearity&2; major synthesis moved to 12.\\
10. Sign tables and inequalities&8.\\
11. Statistics&7, with additions on graphs and the estimated median.\\
12. Lines and systems&3.\\
13. Reference functions (4 parts)&3 for the affine line; 6 for the reference functions and solving; 10 for variations.\\
14. Frequencies and conditional probabilities&11 for frequencies; 14 for conditioning and trees.\\\hline
Progression slots with no direct equivalent&5, 12, 13, 15: reviews, syntheses and consolidation created.
\end{tabularx}\end{center}
The layout is rebuilt in the spirit of the model: A4 portrait, serif typography, black boxes, gridded areas, examples and proofs. The recomposition and the change of order are not a page-by-page facsimile. Pagination is continuous; the contents page no longer uses page numbers that restart at each former sequence.
The usable numerical examples of the model have been kept. Ambiguous figures have been replaced by curves whose expression allows an exact check. The developments on square root, cube, simulation, cross-tabulated data and reviews complete the progression. Added situations with invented data are flagged as fictional; a class data collection is proposed so as to also work on real data.

\section*{Substantive corrections made to the model}
\begin{enumerate}
\item \textbf{Powers.} The repeated product holds for a strictly positive integer exponent; negative exponents require a nonzero base.
\item \textbf{Fractions.} A zero quotient must be defined; only a nonzero fraction has an inverse. Simplifications do not remove the initial restrictions.
\item \textbf{GCD.} The two integers are not both zero in the definition by the greatest divisor; $\pgcd(0,0)=0$ is a separate convention.
\item \textbf{Vectors.} The zero vector directs no line. The characterisation $\vec u=k\vec v$ requires $\vec v\ne\vec0$, or a separate treatment of the zero case. Multiplication by any real number is made explicit.
\item \textbf{Equations of lines.} The condition is $\pt{a;b}\ne\pt{0;0}$; the order of the vectors in the determinant has been made consistent. $\vect{AM}$ can be zero when $M=A$ and is then not a direction vector.
\item \textbf{Intervals.} $0$ is excluded from the set of strictly negative real numbers. Infinities are not real numbers.
\item \textbf{Roots.} The positivity conditions and $\sqrt{x^2}=|x|$ are explicit. A formula that is false in general is not necessarily an inequality that is always true.
\item \textbf{Functions.} Images and preimages are distinguished; the domains of the curves cover the requested values. Graphical readings are not displayed with unjustified precision. The inverse function is not declared decreasing on the whole of $\R^*$.
\item \textbf{Statistics.} The median uses ``at least half''. Adding a value does not make the range decrease. The recalculations use the exact sum $1217$. Variance and standard deviation are not linear functions of the data: their correct transformations are specified.
\item \textbf{Classes.} A continuous data set is not defined by its grouping alone. The mean from the centres and the interpolated median are estimates under assumptions, distinct from the exact values.
\item \textbf{Interpretation.} Multiplying marks is not equated with raising difficulty. Proportion, rate and percentage points are distinguished.
\item \textbf{Probability.} Reference populations and model assumptions are named. Inversion of conditioning, false positives, false negatives and fluctuations are worked on. No general total-probability formula is required.
\item \textbf{Optimisation.} The sweep gives a maximum on a grid; a separate proof establishes the maximum on the interval. A bisection for finding a zero is not presented as a general optimisation method.
\end{enumerate}

\section*{References and coverage check}
\textbf{Progression document used:} \emph{Programmation annuelle Seconde 2026--2027 rectifiée finale}, DOCX supplied. Dates and fifteen stages kept; the dates indicate ranges, not sessions on public holidays.
\par\textbf{Starting document:} \emph{Livret\_2nde\_\_\_Eleve.pdf}, 91 pages, supplied for this adaptation. The original LaTeX source was not attached, so the present source is a reconstruction.
\par\textbf{Official programme:} order of 26 February 2026, BO no.~14 of 2 April 2026, NOR MENE2602914A, applicable from the 2026--2027 school year:
\par\url{https://www.education.gouv.fr/bo/2026/Hebdo14/MENE2602914A}
\par\textbf{Resources supplied:} Seconde programme; resources ``Automatismes 2de'' and ``Raisonnement et démonstration 2de''. The booklet does not claim to reproduce these resources in full.
\begin{center}\small\mxarraystretch{1.5}\begin{tabularx}{\linewidth}{p{8cm} X}
\textbf{Area / skill}&\textbf{Sequences}\\\hline
Integers, multiples, parity, irreducible fractions&1, 5, 8, 13\\
Real numbers, intervals, distance, bounds&4, 5, 8\\
Algebra, powers, roots, equations and signs&8; starter routines all year\\
Vectors, bases, coordinates, norms, determinant&2, 12\\
Equations of lines and intersections&3, 12\\
Representations and six reference functions&6, 10\\
Variations, extrema, comparison, optimisation&10, 13, 15\\
Proportions and percentage change&9, 13, 15\\
Statistics, spread, classes, histogram, cumulative frequencies&7\\
Two-way tables, frequencies, real data&11\\
Conditioning, trees, inversion, simulation&14\\
Logic, proofs, Python, six skills&Spread across the sequences and reviews
\end{tabularx}\end{center}
\textbf{Mathematical check:} the numerical calculations, algebraic expressions, solution sets, table data, coordinates and programs were proofread and checked by calculation. For open collections and productions, the correction gives validity criteria rather than an invented result. Estimates and models remain explicitly flagged.
